% !TEX root = ../main.tex
% Parte feita por todos nós. Combinem no grupo antes de editar este arquivo.

\section{Análise crítica}

\begin{frame}{Problemas de notação}
  \begin{itemize}
    \item Eq. (2): falta o sinal negativo na logística
    \item Eq. (7): falta a raiz e usa desvios-padrão no lugar das variâncias
    \item Eq. (8): faltam os valores 1 e 0 na definição de $\hat{D}$
    \item Eq. (10): $s_i s_i$ e $\delta(c_i, c_i)$ tornam $Q$ independente das comunidades
    \item Tabela 2: KS e Gini trocados para CART e SVM
  \end{itemize}
\end{frame}

\begin{frame}{Limitações do método}
  \begin{itemize}
    \item Nenhum teste de significância para o ganho de AUC
    \item Rede montada com todas as empresas; não explica como pontuar uma empresa nova
    \item MST guarda só $N-1$ arestas
    \item Distância euclidiana sensível a outliers
    \item Acurácia pouco informativa com 11\% de default
  \end{itemize}
\end{frame}

\section{Nossa contribuição}

\begin{frame}{Nossa contribuição}
  \begin{itemize}
    \item \textit{a definir}
  \end{itemize}
\end{frame}

\section{Considerações finais}

\begin{frame}{Referências}
  \small
  \begin{itemize}
    \item Giudici, P., Hadji-Misheva, B. \& Spelta, A. (2019). Network Based Scoring Models to Improve Credit Risk
      Management in Peer to Peer Lending Platforms. \textit{Frontiers in Artificial Intelligence}, 2:3.
    \item Mantegna, R. N. \& Stanley, H. E. (1999). \textit{Introduction to Econophysics}. Cambridge University Press.
    \item Blondel, V. D. et al. (2008). Fast unfolding of communities in large networks. \textit{J. Stat. Mech.}
    \item Newman, M. E. J. (2006). Modularity and community structure in networks. \textit{PNAS}, 103.
  \end{itemize}
\end{frame}

% Slide exigido pelo enunciado (seção 1.6).
\begin{frame}{Ferramentas utilizadas}
  \begin{itemize}
    \item \textit{a preencher}
  \end{itemize}
\end{frame}
